% !TEX program = lualatex
% 情報リテラシー スライドテンプレート

\documentclass[aspectratio=43,professionalfonts,handout]{beamer}
\usepackage{beamer_template}

\newcommand{\LectureNum}{00}          % 回数（01, 02, ...）
\newcommand{\LectureTitle}{タイトル}   % 講義タイトル
\newcommand{\CourseName}{情報リテラシー}
\newcommand{\TextPages}{対応なし}      % 教科書ページ（例: pp.10--20）
\newcommand{\NextTopic}{次回のトピック}
\newcommand{\NextPages}{対応なし}

\begin{document}

% --- Slide 1: タイトル ---

\begin{frame}{\CourseName\quad 第\LectureNum 回「\LectureTitle 」}
  {\small \TermName\quad \TeacherName\quad ／\quad 教科書 \TextPages}

  \vfill
  \begin{block}{今日の到達目標}
    \begin{itemize}
      \item 到達目標をここに書く
    \end{itemize}
  \end{block}

  \begin{block}{チェックリスト}
    \begin{itemize}
      \item[$\square$] チェック項目1
      \item[$\square$] チェック項目2
      \item[$\square$] チェック項目3
    \end{itemize}
  \end{block}
\end{frame}

% --- Slide 2: 基本用語 ---

\begin{frame}{基本用語}
  \begin{description}[labelwidth=6em]
    \item[\kw{用語1}]
      説明文
    \item[\kw{用語2}]
      説明文
    \item[\kw{用語3}]
      説明文
  \end{description}
\end{frame}

% --- Slide 3: ポイント ---

\begin{frame}{ポイント}
  \begin{block}{ポイント}
    \begin{itemize}
      \item ポイント1
      \item ポイント2
      \item ポイント3
    \end{itemize}
  \end{block}
\end{frame}

% --- Slide 4: 演習 ---

\begin{frame}{演習}

  \vfill

  \begin{block}{手順}
    \begin{enumerate}
      \item 手順1
      \item 手順2
      \item 手順3
    \end{enumerate}
  \end{block}
\end{frame}

% --- Slide 5: リンク集（必要に応じて） ---
%
% \begin{frame}{リンク集}
%   \begin{figure}
%     \begin{subfigure}{0.3\linewidth}
%       \includegraphics[width=\linewidth]{images/sample.png}
%       \caption{キャプション}
%     \end{subfigure}
%     \hfill
%     \begin{subfigure}{0.3\linewidth}
%       \includegraphics[width=\linewidth]{images/sample2.png}
%       \caption{キャプション}
%     \end{subfigure}
%     \hfill
%     \begin{subfigure}{0.3\linewidth}
%       \includegraphics[width=\linewidth]{images/sample3.png}
%       \caption{キャプション}
%     \end{subfigure}
%   \end{figure}
% \end{frame}

% --- まとめ ---

\begin{frame}{まとめ}
  \begin{block}{今日のまとめ}
    \begin{itemize}
      \item まとめ1
      \item まとめ2
      \item まとめ3
    \end{itemize}
  \end{block}

  \vfill

  {\small 質問があれば授業後またはTeamsで受け付けます。}
\end{frame}

\end{document}
